Large language model (LLM) unlearning aims to selectively delete specific training data, but unlearning is rarely surgical: removing a forget set $\Df$ also damages knowledge that has nothing to do with $\Df$. Characterising and predicting this collateral damage currently requires running unlearning and an $N\times N$ cross-perplexity (cross-PPL) sweep --- prohibitive at scale. We present two contributions on Llama-3.1-8B-Instruct unlearned with GradAscent across $n{=}100$ forget sets sampled from $10$ HDBSCAN-derived semantic clusters of WikiText-103 ($10{,}000$ unlearned/eval pairs, $1.5\!\times\!10^{6}$ per-sample PPL ratios). \textbf{First}, side-effects decay over three semantic layers but never reach zero: forget samples ($\Lone$) rise to $1.96\times$ baseline PPL, same-cluster neighbours ($\Ltwo$) to $1.32\times$, and cross-cluster bystanders ($\Lthree$) to $1.19\times$, with $73.8\%$ of cross-topic samples still showing elevated PPL. The profile varies across forget sets (L1 spread $1.59\times$ between mildest and worst) --- so \emph{which} forget set you pick matters. \textbf{Second}, we recast benchmarking as a \emph{forget-set audit} problem: a Ridge regressor on $12$-dimensional intrinsic geometric features of $\Df$ predicts the three-layer corruption profile \emph{without ever running unlearning}, attaining held-out Spearman $\rho$ of $0.53$ ($\Lone$), $0.84$ ($\Ltwo$), $0.49$ ($\Lthree$); all three $95\%$ bootstrap CIs exclude zero. The auditor runs in seconds versus several GPU-hours per candidate, enabling a screen-then-run workflow: rank the $100$ candidates by audit score, then unlearn only the top-$k$.
